% =============================================================================
% related_work_section.tex --- Positional Distillation Paper
% Self-contained related work section (~1.5 pages)
% Usage: % =============================================================================
% related_work_section.tex --- Positional Distillation Paper
% Self-contained related work section (~1.5 pages)
% Usage: % =============================================================================
% related_work_section.tex --- Positional Distillation Paper
% Self-contained related work section (~1.5 pages)
% Usage: % =============================================================================
% related_work_section.tex --- Positional Distillation Paper
% Self-contained related work section (~1.5 pages)
% Usage: \input{related_work_section.tex} in main.tex
% =============================================================================

\section{Related Work}
\label{sec:related_work}

\paragraph{Knowledge Distillation for Language Models.}
Knowledge distillation~\citep{hinton2015distilling} transfers knowledge from a teacher model to a smaller student by training on soft targets.
\citet{kim2016sequence} extended this to sequence models with word-level and sequence-level distillation objectives.
For large language models, the choice of divergence measure and data distribution has proven critical.
\citet{gu2024minillm} demonstrated that reverse KL divergence is preferable for generative LLM distillation, as it avoids the student overestimating low-probability teacher modes.
\citet{agarwal2024onpolicy} introduced Generalized Knowledge Distillation (GKD), showing that on-policy generation---where the student generates sequences and the teacher provides soft labels---substantially outperforms supervised (off-policy) distillation.
\citet{ko2024distillm} proposed skew KL divergence and adaptive off-policy approaches for more efficient distillation, while \citet{wen2023fdivergence} generalized sequence-level KD to arbitrary f-divergences.
\citet{wu2025rethinking} challenged the conventional mode-seeking vs.\ mean-seeking interpretation of reverse and forward KL in the LLM setting, finding that both divergences converge given sufficient training but focus on different distribution regions in early epochs.
More recently, \citet{xu2025speculative} introduced Speculative Knowledge Distillation, which uses interleaved teacher-student sampling inspired by speculative decoding, filtering tokens the teacher is unlikely to produce.
Our work builds on the on-policy reverse KL paradigm of \citet{agarwal2024onpolicy} and \citet{gu2024minillm}, but introduces position-based loss computation that fundamentally changes which tokens receive gradient signal.

\paragraph{Token-Level Importance in Distillation and Reasoning.}
A growing body of work demonstrates that not all tokens contribute equally to learning.
\citet{wang2025beyond8020} showed that only $\sim$20\% of chain-of-thought tokens have high entropy, and these ``forking tokens'' steer reasoning paths---applying RL gradients to only these tokens matches full-gradient performance.
\citet{li2025preplan} discovered a ``preplan-and-anchor'' mechanism in LLM attention patterns, where models first perform long-range contextual reference before producing semantic anchor tokens, enabling targeted credit assignment.
\citet{singh2026functional} proposed greedy pruning of reasoning chains and showed that LLMs internally encode token-level functional importance for answer generation.

In the distillation setting specifically, several concurrent works have explored token selection.
\citet{huang2025selectkd} proposed SelecTKD, a propose-and-verify framework where the teacher verifies student-proposed tokens, applying full loss to accepted tokens and masking rejected ones.
\citet{xie2026adakd} introduced AdaKD, which adaptively adjusts token-level temperature based on learning stability.
\citet{tavor2026rethinking} systematically disentangled selective KD along position, class, and sample axes, proposing student-entropy-guided position selection (SE-KD) that reduces wall time by 70\% without sacrificing accuracy.
\citet{kim2026tsdkd} combined entropy-based token selection with preference ranking in a dual distillation framework.

Our approach differs from all of these in a fundamental way: we use \emph{absolute position} as the sole selection criterion, requiring no per-token computation of entropy, difficulty, or verification scores.
This simplicity is motivated by our empirical finding that KL divergence between teacher and student concentrates at early positions corresponding to reasoning strategy selection.
Furthermore, unlike prior token-selection methods that aim primarily at efficiency or robustness, we show that positional loss also eliminates training instability---a qualitative benefit that token-level adaptive methods do not address.

\paragraph{Training Stability and Sequence-Level Issues.}
Full-sequence on-policy distillation is known to suffer from instability.
\citet{holtzman2020curious} identified that likelihood-based training can produce bland, repetitive text---a form of neural text degeneration.
\citet{shumailov2024collapse} demonstrated that iteratively training generative models on their own outputs leads to model collapse, where distribution tails disappear.
In on-policy KD specifically, \citet{jang2026veto} proposed Veto, a geometric bridge in logit space that suppresses harmful gradients on low-confidence tokens, directly addressing the instability we also observe.
Our approach is complementary: rather than modifying the target distribution, we restrict the loss scope to early positions where teacher-student agreement is naturally highest, achieving stability through a simpler mechanism.

Several works have studied how sequence length affects distillation quality.
\citet{chen2025distillingessence} showed that training on only the first 50\% of tokens in supervised reasoning distillation retains $\sim$91\% of full-sequence performance, establishing a computation-quality tradeoff.
\citet{yan2026dasd} addressed exposure bias and capacity misalignment in long chain-of-thought distillation through temperature-scheduled learning.
Our work extends the truncation insight to the on-policy setting: we demonstrate that positional loss masking on student-generated sequences achieves equal or better performance than full-sequence loss, while providing the additional benefit of eliminating degenerate repetition patterns that are unique to on-policy training.

\paragraph{Curriculum Learning and Sequence Scheduling.}
Curriculum learning~\citep{bengio2009curriculum} improves training by gradually increasing task difficulty.
\citet{bengio2015scheduled} applied this principle to sequence models through scheduled sampling, transitioning from teacher-forced to free-running generation during training.
In the LLM pretraining setting, \citet{li2022stability} showed that long sequences at training start cause extreme gradient variance, and sequence length warmup stabilizes training while enabling larger batch sizes.
\citet{pouransari2024dataset} extended this with variable sequence length curriculum for LLM pretraining, achieving up to 3$\times$ faster convergence.
Our progressive position schedule---linearly increasing the position limit from 1 to 200 over training---applies the curriculum principle to the distillation loss scope rather than to the input sequence length, and represents a novel form of curriculum specific to knowledge distillation.

\paragraph{Mathematical Reasoning and Model Context.}
Our experiments use Qwen2.5-Math-1.5B~\citep{qwen2024qwen25math} as the student and Qwen3-1.7B~\citep{qwen2025qwen3} as the teacher, evaluating on the MATH-500 subset~\citep{lightman2023verify} of the MATH benchmark~\citep{hendrycks2021measuring}.
We use LoRA~\citep{hu2022lora} for parameter-efficient fine-tuning.
The distillation-vs-RL paradigm for reasoning has been established by DeepSeek-R1~\citep{deepseekr1}, which showed that distilled models can outperform RL-trained ones in certain settings.
Our work contributes to this landscape by demonstrating that the effectiveness of distillation can be significantly improved through principled control of the loss computation scope.
 in main.tex
% =============================================================================

\section{Related Work}
\label{sec:related_work}

\paragraph{Knowledge Distillation for Language Models.}
Knowledge distillation~\citep{hinton2015distilling} transfers knowledge from a teacher model to a smaller student by training on soft targets.
\citet{kim2016sequence} extended this to sequence models with word-level and sequence-level distillation objectives.
For large language models, the choice of divergence measure and data distribution has proven critical.
\citet{gu2024minillm} demonstrated that reverse KL divergence is preferable for generative LLM distillation, as it avoids the student overestimating low-probability teacher modes.
\citet{agarwal2024onpolicy} introduced Generalized Knowledge Distillation (GKD), showing that on-policy generation---where the student generates sequences and the teacher provides soft labels---substantially outperforms supervised (off-policy) distillation.
\citet{ko2024distillm} proposed skew KL divergence and adaptive off-policy approaches for more efficient distillation, while \citet{wen2023fdivergence} generalized sequence-level KD to arbitrary f-divergences.
\citet{wu2025rethinking} challenged the conventional mode-seeking vs.\ mean-seeking interpretation of reverse and forward KL in the LLM setting, finding that both divergences converge given sufficient training but focus on different distribution regions in early epochs.
More recently, \citet{xu2025speculative} introduced Speculative Knowledge Distillation, which uses interleaved teacher-student sampling inspired by speculative decoding, filtering tokens the teacher is unlikely to produce.
Our work builds on the on-policy reverse KL paradigm of \citet{agarwal2024onpolicy} and \citet{gu2024minillm}, but introduces position-based loss computation that fundamentally changes which tokens receive gradient signal.

\paragraph{Token-Level Importance in Distillation and Reasoning.}
A growing body of work demonstrates that not all tokens contribute equally to learning.
\citet{wang2025beyond8020} showed that only $\sim$20\% of chain-of-thought tokens have high entropy, and these ``forking tokens'' steer reasoning paths---applying RL gradients to only these tokens matches full-gradient performance.
\citet{li2025preplan} discovered a ``preplan-and-anchor'' mechanism in LLM attention patterns, where models first perform long-range contextual reference before producing semantic anchor tokens, enabling targeted credit assignment.
\citet{singh2026functional} proposed greedy pruning of reasoning chains and showed that LLMs internally encode token-level functional importance for answer generation.

In the distillation setting specifically, several concurrent works have explored token selection.
\citet{huang2025selectkd} proposed SelecTKD, a propose-and-verify framework where the teacher verifies student-proposed tokens, applying full loss to accepted tokens and masking rejected ones.
\citet{xie2026adakd} introduced AdaKD, which adaptively adjusts token-level temperature based on learning stability.
\citet{tavor2026rethinking} systematically disentangled selective KD along position, class, and sample axes, proposing student-entropy-guided position selection (SE-KD) that reduces wall time by 70\% without sacrificing accuracy.
\citet{kim2026tsdkd} combined entropy-based token selection with preference ranking in a dual distillation framework.

Our approach differs from all of these in a fundamental way: we use \emph{absolute position} as the sole selection criterion, requiring no per-token computation of entropy, difficulty, or verification scores.
This simplicity is motivated by our empirical finding that KL divergence between teacher and student concentrates at early positions corresponding to reasoning strategy selection.
Furthermore, unlike prior token-selection methods that aim primarily at efficiency or robustness, we show that positional loss also eliminates training instability---a qualitative benefit that token-level adaptive methods do not address.

\paragraph{Training Stability and Sequence-Level Issues.}
Full-sequence on-policy distillation is known to suffer from instability.
\citet{holtzman2020curious} identified that likelihood-based training can produce bland, repetitive text---a form of neural text degeneration.
\citet{shumailov2024collapse} demonstrated that iteratively training generative models on their own outputs leads to model collapse, where distribution tails disappear.
In on-policy KD specifically, \citet{jang2026veto} proposed Veto, a geometric bridge in logit space that suppresses harmful gradients on low-confidence tokens, directly addressing the instability we also observe.
Our approach is complementary: rather than modifying the target distribution, we restrict the loss scope to early positions where teacher-student agreement is naturally highest, achieving stability through a simpler mechanism.

Several works have studied how sequence length affects distillation quality.
\citet{chen2025distillingessence} showed that training on only the first 50\% of tokens in supervised reasoning distillation retains $\sim$91\% of full-sequence performance, establishing a computation-quality tradeoff.
\citet{yan2026dasd} addressed exposure bias and capacity misalignment in long chain-of-thought distillation through temperature-scheduled learning.
Our work extends the truncation insight to the on-policy setting: we demonstrate that positional loss masking on student-generated sequences achieves equal or better performance than full-sequence loss, while providing the additional benefit of eliminating degenerate repetition patterns that are unique to on-policy training.

\paragraph{Curriculum Learning and Sequence Scheduling.}
Curriculum learning~\citep{bengio2009curriculum} improves training by gradually increasing task difficulty.
\citet{bengio2015scheduled} applied this principle to sequence models through scheduled sampling, transitioning from teacher-forced to free-running generation during training.
In the LLM pretraining setting, \citet{li2022stability} showed that long sequences at training start cause extreme gradient variance, and sequence length warmup stabilizes training while enabling larger batch sizes.
\citet{pouransari2024dataset} extended this with variable sequence length curriculum for LLM pretraining, achieving up to 3$\times$ faster convergence.
Our progressive position schedule---linearly increasing the position limit from 1 to 200 over training---applies the curriculum principle to the distillation loss scope rather than to the input sequence length, and represents a novel form of curriculum specific to knowledge distillation.

\paragraph{Mathematical Reasoning and Model Context.}
Our experiments use Qwen2.5-Math-1.5B~\citep{qwen2024qwen25math} as the student and Qwen3-1.7B~\citep{qwen2025qwen3} as the teacher, evaluating on the MATH-500 subset~\citep{lightman2023verify} of the MATH benchmark~\citep{hendrycks2021measuring}.
We use LoRA~\citep{hu2022lora} for parameter-efficient fine-tuning.
The distillation-vs-RL paradigm for reasoning has been established by DeepSeek-R1~\citep{deepseekr1}, which showed that distilled models can outperform RL-trained ones in certain settings.
Our work contributes to this landscape by demonstrating that the effectiveness of distillation can be significantly improved through principled control of the loss computation scope.
 in main.tex
% =============================================================================

\section{Related Work}
\label{sec:related_work}

\paragraph{Knowledge Distillation for Language Models.}
Knowledge distillation~\citep{hinton2015distilling} transfers knowledge from a teacher model to a smaller student by training on soft targets.
\citet{kim2016sequence} extended this to sequence models with word-level and sequence-level distillation objectives.
For large language models, the choice of divergence measure and data distribution has proven critical.
\citet{gu2024minillm} demonstrated that reverse KL divergence is preferable for generative LLM distillation, as it avoids the student overestimating low-probability teacher modes.
\citet{agarwal2024onpolicy} introduced Generalized Knowledge Distillation (GKD), showing that on-policy generation---where the student generates sequences and the teacher provides soft labels---substantially outperforms supervised (off-policy) distillation.
\citet{ko2024distillm} proposed skew KL divergence and adaptive off-policy approaches for more efficient distillation, while \citet{wen2023fdivergence} generalized sequence-level KD to arbitrary f-divergences.
\citet{wu2025rethinking} challenged the conventional mode-seeking vs.\ mean-seeking interpretation of reverse and forward KL in the LLM setting, finding that both divergences converge given sufficient training but focus on different distribution regions in early epochs.
More recently, \citet{xu2025speculative} introduced Speculative Knowledge Distillation, which uses interleaved teacher-student sampling inspired by speculative decoding, filtering tokens the teacher is unlikely to produce.
Our work builds on the on-policy reverse KL paradigm of \citet{agarwal2024onpolicy} and \citet{gu2024minillm}, but introduces position-based loss computation that fundamentally changes which tokens receive gradient signal.

\paragraph{Token-Level Importance in Distillation and Reasoning.}
A growing body of work demonstrates that not all tokens contribute equally to learning.
\citet{wang2025beyond8020} showed that only $\sim$20\% of chain-of-thought tokens have high entropy, and these ``forking tokens'' steer reasoning paths---applying RL gradients to only these tokens matches full-gradient performance.
\citet{li2025preplan} discovered a ``preplan-and-anchor'' mechanism in LLM attention patterns, where models first perform long-range contextual reference before producing semantic anchor tokens, enabling targeted credit assignment.
\citet{singh2026functional} proposed greedy pruning of reasoning chains and showed that LLMs internally encode token-level functional importance for answer generation.

In the distillation setting specifically, several concurrent works have explored token selection.
\citet{huang2025selectkd} proposed SelecTKD, a propose-and-verify framework where the teacher verifies student-proposed tokens, applying full loss to accepted tokens and masking rejected ones.
\citet{xie2026adakd} introduced AdaKD, which adaptively adjusts token-level temperature based on learning stability.
\citet{tavor2026rethinking} systematically disentangled selective KD along position, class, and sample axes, proposing student-entropy-guided position selection (SE-KD) that reduces wall time by 70\% without sacrificing accuracy.
\citet{kim2026tsdkd} combined entropy-based token selection with preference ranking in a dual distillation framework.

Our approach differs from all of these in a fundamental way: we use \emph{absolute position} as the sole selection criterion, requiring no per-token computation of entropy, difficulty, or verification scores.
This simplicity is motivated by our empirical finding that KL divergence between teacher and student concentrates at early positions corresponding to reasoning strategy selection.
Furthermore, unlike prior token-selection methods that aim primarily at efficiency or robustness, we show that positional loss also eliminates training instability---a qualitative benefit that token-level adaptive methods do not address.

\paragraph{Training Stability and Sequence-Level Issues.}
Full-sequence on-policy distillation is known to suffer from instability.
\citet{holtzman2020curious} identified that likelihood-based training can produce bland, repetitive text---a form of neural text degeneration.
\citet{shumailov2024collapse} demonstrated that iteratively training generative models on their own outputs leads to model collapse, where distribution tails disappear.
In on-policy KD specifically, \citet{jang2026veto} proposed Veto, a geometric bridge in logit space that suppresses harmful gradients on low-confidence tokens, directly addressing the instability we also observe.
Our approach is complementary: rather than modifying the target distribution, we restrict the loss scope to early positions where teacher-student agreement is naturally highest, achieving stability through a simpler mechanism.

Several works have studied how sequence length affects distillation quality.
\citet{chen2025distillingessence} showed that training on only the first 50\% of tokens in supervised reasoning distillation retains $\sim$91\% of full-sequence performance, establishing a computation-quality tradeoff.
\citet{yan2026dasd} addressed exposure bias and capacity misalignment in long chain-of-thought distillation through temperature-scheduled learning.
Our work extends the truncation insight to the on-policy setting: we demonstrate that positional loss masking on student-generated sequences achieves equal or better performance than full-sequence loss, while providing the additional benefit of eliminating degenerate repetition patterns that are unique to on-policy training.

\paragraph{Curriculum Learning and Sequence Scheduling.}
Curriculum learning~\citep{bengio2009curriculum} improves training by gradually increasing task difficulty.
\citet{bengio2015scheduled} applied this principle to sequence models through scheduled sampling, transitioning from teacher-forced to free-running generation during training.
In the LLM pretraining setting, \citet{li2022stability} showed that long sequences at training start cause extreme gradient variance, and sequence length warmup stabilizes training while enabling larger batch sizes.
\citet{pouransari2024dataset} extended this with variable sequence length curriculum for LLM pretraining, achieving up to 3$\times$ faster convergence.
Our progressive position schedule---linearly increasing the position limit from 1 to 200 over training---applies the curriculum principle to the distillation loss scope rather than to the input sequence length, and represents a novel form of curriculum specific to knowledge distillation.

\paragraph{Mathematical Reasoning and Model Context.}
Our experiments use Qwen2.5-Math-1.5B~\citep{qwen2024qwen25math} as the student and Qwen3-1.7B~\citep{qwen2025qwen3} as the teacher, evaluating on the MATH-500 subset~\citep{lightman2023verify} of the MATH benchmark~\citep{hendrycks2021measuring}.
We use LoRA~\citep{hu2022lora} for parameter-efficient fine-tuning.
The distillation-vs-RL paradigm for reasoning has been established by DeepSeek-R1~\citep{deepseekr1}, which showed that distilled models can outperform RL-trained ones in certain settings.
Our work contributes to this landscape by demonstrating that the effectiveness of distillation can be significantly improved through principled control of the loss computation scope.
 in main.tex
% =============================================================================

\section{Related Work}
\label{sec:related_work}

\paragraph{Knowledge Distillation for Language Models.}
Knowledge distillation~\citep{hinton2015distilling} transfers knowledge from a teacher model to a smaller student by training on soft targets.
\citet{kim2016sequence} extended this to sequence models with word-level and sequence-level distillation objectives.
For large language models, the choice of divergence measure and data distribution has proven critical.
\citet{gu2024minillm} demonstrated that reverse KL divergence is preferable for generative LLM distillation, as it avoids the student overestimating low-probability teacher modes.
\citet{agarwal2024onpolicy} introduced Generalized Knowledge Distillation (GKD), showing that on-policy generation---where the student generates sequences and the teacher provides soft labels---substantially outperforms supervised (off-policy) distillation.
\citet{ko2024distillm} proposed skew KL divergence and adaptive off-policy approaches for more efficient distillation, while \citet{wen2023fdivergence} generalized sequence-level KD to arbitrary f-divergences.
\citet{wu2025rethinking} challenged the conventional mode-seeking vs.\ mean-seeking interpretation of reverse and forward KL in the LLM setting, finding that both divergences converge given sufficient training but focus on different distribution regions in early epochs.
More recently, \citet{xu2025speculative} introduced Speculative Knowledge Distillation, which uses interleaved teacher-student sampling inspired by speculative decoding, filtering tokens the teacher is unlikely to produce.
Our work builds on the on-policy reverse KL paradigm of \citet{agarwal2024onpolicy} and \citet{gu2024minillm}, but introduces position-based loss computation that fundamentally changes which tokens receive gradient signal.

\paragraph{Token-Level Importance in Distillation and Reasoning.}
A growing body of work demonstrates that not all tokens contribute equally to learning.
\citet{wang2025beyond8020} showed that only $\sim$20\% of chain-of-thought tokens have high entropy, and these ``forking tokens'' steer reasoning paths---applying RL gradients to only these tokens matches full-gradient performance.
\citet{li2025preplan} discovered a ``preplan-and-anchor'' mechanism in LLM attention patterns, where models first perform long-range contextual reference before producing semantic anchor tokens, enabling targeted credit assignment.
\citet{singh2026functional} proposed greedy pruning of reasoning chains and showed that LLMs internally encode token-level functional importance for answer generation.

In the distillation setting specifically, several concurrent works have explored token selection.
\citet{huang2025selectkd} proposed SelecTKD, a propose-and-verify framework where the teacher verifies student-proposed tokens, applying full loss to accepted tokens and masking rejected ones.
\citet{xie2026adakd} introduced AdaKD, which adaptively adjusts token-level temperature based on learning stability.
\citet{tavor2026rethinking} systematically disentangled selective KD along position, class, and sample axes, proposing student-entropy-guided position selection (SE-KD) that reduces wall time by 70\% without sacrificing accuracy.
\citet{kim2026tsdkd} combined entropy-based token selection with preference ranking in a dual distillation framework.

Our approach differs from all of these in a fundamental way: we use \emph{absolute position} as the sole selection criterion, requiring no per-token computation of entropy, difficulty, or verification scores.
This simplicity is motivated by our empirical finding that KL divergence between teacher and student concentrates at early positions corresponding to reasoning strategy selection.
Furthermore, unlike prior token-selection methods that aim primarily at efficiency or robustness, we show that positional loss also eliminates training instability---a qualitative benefit that token-level adaptive methods do not address.

\paragraph{Training Stability and Sequence-Level Issues.}
Full-sequence on-policy distillation is known to suffer from instability.
\citet{holtzman2020curious} identified that likelihood-based training can produce bland, repetitive text---a form of neural text degeneration.
\citet{shumailov2024collapse} demonstrated that iteratively training generative models on their own outputs leads to model collapse, where distribution tails disappear.
In on-policy KD specifically, \citet{jang2026veto} proposed Veto, a geometric bridge in logit space that suppresses harmful gradients on low-confidence tokens, directly addressing the instability we also observe.
Our approach is complementary: rather than modifying the target distribution, we restrict the loss scope to early positions where teacher-student agreement is naturally highest, achieving stability through a simpler mechanism.

Several works have studied how sequence length affects distillation quality.
\citet{chen2025distillingessence} showed that training on only the first 50\% of tokens in supervised reasoning distillation retains $\sim$91\% of full-sequence performance, establishing a computation-quality tradeoff.
\citet{yan2026dasd} addressed exposure bias and capacity misalignment in long chain-of-thought distillation through temperature-scheduled learning.
Our work extends the truncation insight to the on-policy setting: we demonstrate that positional loss masking on student-generated sequences achieves equal or better performance than full-sequence loss, while providing the additional benefit of eliminating degenerate repetition patterns that are unique to on-policy training.

\paragraph{Curriculum Learning and Sequence Scheduling.}
Curriculum learning~\citep{bengio2009curriculum} improves training by gradually increasing task difficulty.
\citet{bengio2015scheduled} applied this principle to sequence models through scheduled sampling, transitioning from teacher-forced to free-running generation during training.
In the LLM pretraining setting, \citet{li2022stability} showed that long sequences at training start cause extreme gradient variance, and sequence length warmup stabilizes training while enabling larger batch sizes.
\citet{pouransari2024dataset} extended this with variable sequence length curriculum for LLM pretraining, achieving up to 3$\times$ faster convergence.
Our progressive position schedule---linearly increasing the position limit from 1 to 200 over training---applies the curriculum principle to the distillation loss scope rather than to the input sequence length, and represents a novel form of curriculum specific to knowledge distillation.

\paragraph{Mathematical Reasoning and Model Context.}
Our experiments use Qwen2.5-Math-1.5B~\citep{qwen2024qwen25math} as the student and Qwen3-1.7B~\citep{qwen2025qwen3} as the teacher, evaluating on the MATH-500 subset~\citep{lightman2023verify} of the MATH benchmark~\citep{hendrycks2021measuring}.
We use LoRA~\citep{hu2022lora} for parameter-efficient fine-tuning.
The distillation-vs-RL paradigm for reasoning has been established by DeepSeek-R1~\citep{deepseekr1}, which showed that distilled models can outperform RL-trained ones in certain settings.
Our work contributes to this landscape by demonstrating that the effectiveness of distillation can be significantly improved through principled control of the loss computation scope.
